% Part: first-order-logic
% Chapter: completeness
% Section: downward-ls

\documentclass[../../../include/open-logic-section]{subfiles}

\begin{document}

\olfileid{fol}{com}{dls}
\olsection{مبرهنة لوفنهايم--سكولم}

مقتضى مبرهنة لوفنهايم--سكولم أن النظرية إذا كان لها نموذج غير منتهٍ،
كان لها نموذج !!{denumerable} على الأكثر. فيلزم من ذلك عجز منطق
الرتبة الأولى عن التعبير عن أن حجم !!a{structure} غير
!!{enumerable}: فكل !!{sentence}، أو مجموعة من ال!!{sentence}s،
تصدق في جميع ال!!{structure}s غير ال!!{enumerable} لا بد أن
تصدق في بنية !!{enumerable} أيضًا.

\begin{thm}
\ollabel{thm:downward-ls} إذا كانت $\Gamma$ متسقة، فلها نموذج
!!a{enumerable}، أي إنها قابلة للإرضاء في !!a{structure} مجالها إما
منتهٍ وإما !!{denumerable}.
\end{thm}

\begin{proof}
لتكن $\Gamma$ متسقة في~$\Lang L$. برهان الاكتمال يمد اللغة
امتدادًا هنكنيًا~$\Lang{L'}$، ثم يبني بنية حدود لا يتجاوز حجم
مجالها حجم مجموعة الحدود المغلقة~$\Trm[L']_0$ في~$\Lang{L'}$.
فهي !!{denumerable} على الأكثر. وإذا رددناها~$\Struct M$ إلى
اللغة الأصلية~$\Lang L$، حصل نموذج لـ~$\Gamma$ لم يتغير مجاله.
\end{proof}

\begin{thm}
\ollabel{noidentity-ls} إذا كانت $\Gamma$ مجموعة متسقة من ال!!{sentence}s
في لغة منطق الرتبة الأولى الخالية من ال!!{identity}، فلها نموذج
!!a{denumerable}، أي إنها قابلة للإرضاء في !!a{structure} مجالها غير
منتهٍ و!!{enumerable}.
\end{thm}

\begin{proof}
لتكن $\Gamma$ متسقة، ولا يقع في أي !!{sentence} منها رمز
ال!!{identity}. حينئذ تكون بنية الحدود المنشأة في لغة
هنكن~$\Lang L'$ ببرهان الاكتمال ذات مجال~$\Domain{M}$ هو بعينه
مجموعة الحدود المغلقة لتلك اللغة. فرَدُّ~$\Struct{M}$ إلى اللغة
الأصلية نموذج !!{denumerable} لـ~$\Gamma$، لأن
$\Trm[L']_0$ !!{denumerable}.
\end{proof}

\begin{ex}[مفارقة سكولم]
تتيح نظرية مجموعات زرميلو--فرينكل~$\Log{ZFC}$ التعبير عن العبارات
الرياضية كلها تقريبًا، ومن جملتها أحجام المجموعات. فمن مبرهنات
$\Log{ZFC}$ أن مجموعة الأعداد الحقيقية~$\Real$ غير
!!{enumerable}، ومبرهنة كانتور القاضية بأن مجموعة الأجزاء أكبر من
المجموعة الأصلية، وغير ذلك. فإذا كانت $\Log{ZFC}$ متسقة، لم يكن
لها إلا نماذج غير منتهية؛ وفي كل منها !!{element}s تحكم النظرية
بأنها غير !!{enumerable}. ومن أمثلتها العنصر الشاهد لمبرهنة
~$\Log{ZFC}$ بوجود مجموعة أجزاء الأعداد الطبيعية. ومع ذلك تعطي
مبرهنة لوفنهايم--سكولم لـ~$\Log{ZFC}$ نماذج !!{enumerable}.
ففي هذه النماذج مجموعات «غير !!{enumerable}»، والنماذج نفسها
!!{enumerable}.
\end{ex}

\end{document}
